\chapter{IMPLEMENTATION AND TESTING}

\section{Implementation}

\subsection{Tools Used}
The project was implemented using the following set of tools and technologies:

\begin{itemize}
    \item \textbf{Frontend:} React.js (v18), TailwindCSS (Styling), Recharts (Data Visualization), Capacitor (Mobile Build).
    \item \textbf{Backend:} Python 3.10+, FastAPI (Web Framework), Uvicorn (ASGI Server), Google Gemini SDK (AI).
    \item \textbf{Database Platform:} Supabase (PostgreSQL), Supabase Auth.
    \item \textbf{CASE Tools:} Visual Studio Code, Git, GitHub.
\end{itemize}

\subsection{Implementation Details of Modules}

\subsubsection{User Authentication Module}
Implemented using Supabase Auth. Provides functions for \texttt{signUp()}, \texttt{signInWithPassword()}, and \texttt{signOut()}. Handles session management and JWT token verification.

\subsubsection{Expense Parsing Module (Backend)}
The core feature of the application is the NLP parser. Below is the Python implementation using FastAPI and Gemini:

\begin{lstlisting}[language=Python, caption=Backend API for Expense Parsing]
# backend/api/expenses.py (Snippet)

@router.post("/parse")
async def parse_expense(text: str):
    """
    Parses natural language text into structured expense data.
    """
    prompt = f"""
    Extract the following details from the text: "{text}"
    - Amount (number)
    - Category (Enum: Food, Travel, Bills, Entertainment, Health, Other)
    - Description (string)
    
    Return JSON format only.
    """
    
    response = model.generate_content(prompt)
    try:
        data = json.loads(response.text)
        return {"status": "success", "data": data}
    except json.JSONDecodeError:
        return {"status": "error", "message": "Failed to parse"}
\end{lstlisting}

\subsubsection{Expense Management Module (Frontend)}
The frontend manages state using React Context API to ensure data consistency across components.

\begin{lstlisting}[language=JavaScript, caption=React Context for Expenses]
// src/context/ExpenseContext.js (Snippet)

export const ExpenseProvider = ({ children }) => {
    const [expenses, setExpenses] = useState([]);

    const addExpense = async (text) => {
        const parsed = await api.parseExpense(text);
        if (parsed.status === "success") {
            const newExpense = await api.saveExpense(parsed.data);
            setExpenses([newExpense, ...expenses]);
        }
    };

    return (
        <ExpenseContext.Provider value={{ expenses, addExpense }}>
            {children}
        </ExpenseContext.Provider>
    );
};
\end{lstlisting}

\section{Testing}

\subsection{Test Cases for Unit Testing}
Unit tests were written for individual backend functions, particularly the NLP parsing logic.

\begin{table}[H]
\centering
\begin{tabular}{|l|l|l|l|}
\hline
\textbf{Test ID} & \textbf{Input} & \textbf{Expected Output} & \textbf{Result} \\ \hline
UT-01 & "Food 500" & \{amount: 500, category: "Food"\} & Pass \\ \hline
UT-02 & "Taxi 200" & \{amount: 200, category: "Travel"\} & Pass \\ \hline
UT-03 & "Invalid" & Error Message & Pass \\ \hline
\end{tabular}
\caption{Unit Test Cases}
\end{table}

\subsection{Test Cases for System Testing}
System tests verified the end-to-end flow.

\begin{table}[H]
\centering
\begin{tabular}{|l|l|l|l|}
\hline
\textbf{Test ID} & \textbf{Scenario} & \textbf{Expected Outcome} & \textbf{Result} \\ \hline
ST-01 & User Logs In & Dashboard loads with user data & Pass \\ \hline
ST-02 & Add Expense via NLP & Expense appears in list & Pass \\ \hline
ST-03 & View Reports & Charts render with correct data & Pass \\ \hline
ST-04 & Mobile View & Layout adapts to screen size & Pass \\ \hline
\end{tabular}
\caption{System Test Cases}
\end{table}
